\chapter{Conclusão}\label{cap:conclusao}

Investigaram-se a coerência \textit{Wavelet} alternativa e a semelhança dos perfis
espectrais como descritores de padrões de escoamento bifásico em microtubo horizontal com
isobutano. A análise utilizou sinais de intensidade extraídos das linhas 18 a 71 dos
panoramas de 26 ensaios, distribuídos em cinco padrões. Foram examinados 325 pares de
ensaios distintos: 87 do mesmo padrão e 238 de padrões diferentes.

A comparação da forma espectral apresentou um resultado promissor. O coeficiente de
Pearson $r$, calculado entre os resíduos dos perfis logarítmicos de potência no atraso
zero, teve média de $0{,}8697$ dentro dos padrões e de $0{,}8038$ entre padrões. A diferença
positiva de $0{,}0659$ indicou maior semelhança média no primeiro grupo. O padrão
\textit{plug} e o anular ondulado apresentaram as maiores médias internas do coeficiente
de Pearson $r$, respectivamente $0{,}9618$ e $0{,}9045$. Entretanto, também ocorreram
correlações elevadas entre padrões distintos, com ampla sobreposição das distribuições.
Esses resultados sustentaram o aprofundamento da análise dos perfis, sem estabelecer uma
regra de classificação ou demonstrar significância estatística nesta etapa.

O escore ponderado de coerência apresentou média de $0{,}1877$ dentro dos padrões e de
$0{,}2635$ entre padrões. Assim, esse resumo não reproduziu o comportamento de maior
semelhança interna pretendido para a discriminação global. A diferença em relação ao
coeficiente de Pearson $r$ foi interpretada considerando as definições de cada medida:
o coeficiente de Pearson $r$ comparou a forma dos perfis dos panoramas completos, enquanto
o escore de coerência reuniu máximos por escala nos prefixos comuns de cada par. Os máximos
puderam ocorrer em atrasos diferentes, sem corresponder a um único alinhamento físico.

A conferência das entradas acrescentou uma condição importante à interpretação. No ensaio
anular \codigo{m03q40\_001}, identificaram-se 60 substituições de deslocamento em 500
intervalos e $25{,}43\%$ dos pixels da faixa analisada no teto de intensidade. Na análise
exploratória sem esse registro, a média interna do coeficiente de Pearson $r$ do padrão
anular aumentou para $0{,}9508$, mas a diferença global entre as médias intra-padrão e
inter-padrão diminuiu para $0{,}0549$. Portanto, a tendência favorável dos perfis persistiu,
sem que a retirada demonstrasse melhora global da discriminação. O conjunto original foi
mantido como referência, acompanhado dessa sensibilidade.

O objetivo de investigar os descritores foi atendido pela aplicação das equações, pela
comparação dos dois tipos de medida e pela documentação das condições em que os resultados
foram obtidos. A revisão bibliográfica e a fundamentação apoiaram a metodologia; as figuras
e tabelas permitiram acompanhar a passagem dos panoramas aos sinais, aos espectros e aos
resumos. A verificação numérica confirmou as operações implementadas, enquanto a análise
de qualidade explicitou limitações da aquisição e da reconstrução que permaneceram
relevantes para a interpretação física.

As evidências obtidas constituíram uma base para novas análises. Permaneceram em aberto a
avaliação inferencial com o ensaio como unidade de replicação, a influência das condições
operacionais e ópticas e a definição de critérios gerais de qualidade da reconstrução.
Identificaram-se como desdobramentos a ampliação dos ensaios, especialmente para os padrões
com poucas aquisições, o exame de descritores complementares e sua posterior avaliação em
dados separados por aquisição. Esses encaminhamentos foram indicados para investigar a
reprodutibilidade da discriminação, preservando a distinção entre o indício descritivo já
obtido e a validação ainda necessária.
